\scanpage{26}
% L1–2
\diagram{\begin{tikzcd}[ampersand replacement=\&]
p^{(n)}(Y/X/S)\arrow[rr,"{\unsure{\varepsilon|E}}^n(Y/X/S)"]\arrow[dr]\&\&Y\arrow[dl]\\
\&X\&
\end{tikzcd}}
% L2–4
d'autre part, on va définir un \uline{hom} \deleted{application} \uline{de faisceaux} sur $X$
\[
d^n(Y/X/S):\underline{\Gamma}(Y/X)\longrightarrow\underline{\Gamma}(P^{(n)}(Y/X/S)/X)
\]
% slanted
\margin{\uline{\deleted{\unsure{système de}}} \uline{parties principales}}
% L5–8
[tel que le composé avec $\Gamma$
\[
{\unsure{g|D|\delta|\zeta|\varepsilon}}^n(Y/X/S):\underline{\Gamma}(p^{(n)}(Y/X/S)/X)\longrightarrow\underline{\Gamma}(Y/X)
\]
soit l'identité :
\[
g^n d^n=\operatorname{id}\qquad\big]
\]

% L8, small isolated mark below display, at left
\margin{$\unsure{x|\times}$}
% L9–12
\deleted{Il faut donc}

Pour ceci, \deleted{\unsure{on note que}} \inserted{il suffit de définir}
\[
\Gamma(Y/X)\longrightarrow\Gamma(P^n(Y/X/S)/X)
\]
et de vérifier sa nature locale sur $X$. On note que
% L13–14
\[
\Gamma(Y/X)\longrightarrow\deleted{\operatorname{Hom}}\Gamma(Y\times_X\Delta^{(n)}_{X/S}/\Delta^{(n)}_{X/S})
\]
% L15–20
Or
\[
\begin{aligned}
\operatorname{Hom}_X(\Delta^n_{X/S},Y)&=\Gamma(Y\times_X\Delta^{(n)}_{X/S}/\Delta^{(n)}_{X/S})\simeq\text{\unsure{schéma|ensemble|réunion} des}\\
&\text{sections d'\unsure{ordre|ord.|au-dessus}\deleted{\unsure{x}} $n$ de }Y\times_S X\text{ sur}\\
&X\times_S X\text{ le long de l'immersion}\\
&X\longrightarrow X\times_S X
\end{aligned}
\]
et par \unsure{symétrie|hypothèse|projections}, ce dernier ensemble est
% archive sheet number: 25

